\section{A variational principle for tangent vector fields}\label{sec:vector}

For a real vector field $X\in H^1(\Omega;\R^2)$ with tangent trace, meaning $\tr X\cdot\nu=0$, set
\[
\Q(X)=\int_\Omega\bigl(|\diver X|^2+|\curl X|^2\bigr),
\qquad \curl X=\partial_1X_2-\partial_2X_1.
\]
The following is Rohleder's vector variational principle \cite[Theorem~1.1]{Rohleder2021}; see also \cite{Rohleder2024}. We include its derivation, in particular its equality case. Our curvature convention below is the opposite of his, so the corresponding boundary term has the opposite sign.

\begin{proposition}\label{prop:vector}
For every such $X$,
\begin{equation}\label{eq:vector-inequality}
\Q(X)\ge\mu\int_\Omega|X|^2,
\end{equation}
with equality exactly when $X=\grad v$ for some $v\in\V$.
\end{proposition}
\begin{proof}
Solve
\[
\Delta v=\diver X,\quad \partial_\nu v=0,\quad\int_\Omega v=0,
\qquad
\Delta j=\curl X,\quad j|_{\partial\Omega}=0.
\]
The Neumann compatibility follows from the tangent trace. The solutions are in $H^2(\Omega)$ by smooth-boundary regularity \cite[Chapter~III, Theorems~6.1 and~6.4]{Showalter1994}. The rotated gradient $\grad^\perp j=(-\partial_2j,\partial_1j)$ is tangent. Thus
\[
Z=X-\grad v-\grad^\perp j
\]
is tangent, curl-free, and divergence-free. Its components are harmonic in the interior, so $Z$ is smooth there. Simple connectivity gives an interior potential $r$ with $Z=\grad r$.

Here $r\in H^1(\Omega)$, as can be seen without assuming its global integrability in advance. For $M>0$, its bounded truncation $r_M=\max(-M,\min(r,M))$ has local weak gradient of norm at most $|Z|$, hence belongs globally to $H^1(\Omega)$. Fix an interior ball $B$. Its averages $(r_M)_B$ are bounded by $\sup_B|r|$. The Poincar\'e inequality, with this fixed-ball average as anchor, gives
\[
\norm{r_M}_{L^2(\Omega)}
\le C\norm{\grad r_M}_{L^2(\Omega)}+C|(r_M)_B|
\le C\norm Z_2+C\sup_B|r|.
\]
Fatou gives $r\in L^2$, and its weak gradient is $Z$. Testing $\diver Z=0$ against $r$, with its zero normal trace, yields $\int|\grad r|^2=0$. Thus
\[
X=\grad v+\grad^\perp j.
\]
The summands are orthogonal in $L^2$, since the second is divergence-free and tangent. Moreover,
\[
\Q(X)=\norm{\Delta v}_2^2+\norm{\Delta j}_2^2.
\]
Integration by parts and the two Poincar\'e inequalities show
\[
\norm{\grad v}_2^2
\le\norm v_2\norm{\Delta v}_2
\le\mu^{-1/2}\norm{\grad v}_2\norm{\Delta v}_2,
\]
and likewise $\norm{\Delta j}_2^2\ge\lambda_D\norm{\grad j}_2^2$. Hence
\[
\Q(X)\ge\mu\norm{\grad v}_2^2+\lambda_D\norm{\grad j}_2^2
=\mu\norm X_2^2+(\lambda_D-\mu)\norm{\grad j}_2^2.
\]
By \Cref{lem:gap}, this proves \eqref{eq:vector-inequality}. Equality forces $j=0$ and equality in the Rayleigh inequality for $v$, so $v\in\V$. Conversely, gradients of functions in $\V$ give equality directly.
\end{proof}

For a differential-form interpretation of this divergence--curl energy, see \cite[Section~5.1]{FriesGoffengMiranda2026}. We will use the following boundary identity in the fixed Cartesian coordinates of the plane.

Let $\tau$ be the counterclockwise unit tangent to $\partial\Omega$, and put
\[
\kappa=\det(\tau,\partial_\tau\tau).
\]
Thus $\kappa$ is positive on a convex circle.

\begin{lemma}\label{lem:curvature}
For tangent $X\in H^1(\Omega;\R^2)$,
\begin{equation}\label{eq:curvature}
\Q(X)=\int_\Omega|\grad X|^2
+\int_{\partial\Omega}\kappa|X|^2.
\end{equation}
\end{lemma}
\begin{proof}
For smooth fields, the difference of the bulk integrands is $2\det\grad X$. Its integral is the boundary pairing
\[
\int_{\partial\Omega}
(X_1\partial_\tau X_2-X_2\partial_\tau X_1).
\]
This identity extends to $H^1$ fields by smooth approximation, since traces are in $H^{1/2}$ and their tangential derivatives are in $H^{-1/2}$. If the trace is $a\tau$, the pairing is $\int\kappa a^2$. This follows first for smooth $a$ and then by approximation in $H^{1/2}$. It gives \eqref{eq:curvature}.
\end{proof}
